% © 2026 Ing. David Jaroš — CC BY-NC-ND 4.0
% UBT-AI-PROVENANCE-BEGIN
% schema: ubt-ai-provenance/v1
% tier: B_machine_verified
% ai_assistance: disclosed
% human_review: machine-verification
% editorial_responsibility: Ing. David Jaroš
% policy: ../../AI_PROVENANCE.md
% notice: Exact algebraic no-go claims are machine-verified; dynamical claims remain conditional/open.
% UBT-AI-PROVENANCE-END

\documentclass[11pt,a4paper]{article}
\usepackage[margin=2.5cm]{geometry}
\usepackage{amsmath,amssymb,amsthm,mathtools}
\usepackage[most]{tcolorbox}
\newtheorem{theorem}{Theorem}
\newtheorem{corollary}[theorem]{Corollary}
\newtheorem{proposition}[theorem]{Proposition}
\newtheorem{remark}[theorem]{Remark}
\newcommand{\C}{\mathbb C}
\newcommand{\Hq}{\mathbb H}
\newcommand{\Aalg}{\mathcal A}
\newcommand{\End}{\operatorname{End}}
\newcommand{\Sc}{\operatorname{Sc}}
\title{The SU(3) Dynamical Boundary in UBT:\\
A Minimal-Bimodule No-Go and a Conditional Induced Yang--Mills Route}
\author{Ing.\ David Jaro\v{s}}
\date{9 October 2026}
\begin{document}
\maketitle

\begin{abstract}
The canonical colour carrier
$V=\C\operatorname{-span}\{I,J,K\}$ has the algebraically established
stabiliser $\operatorname{Stab}(h,\Omega)=SU(3)$.  This note determines what
does and does not follow dynamically from the minimal two-sided
biquaternionic derivative.  We prove that every single operator
$X\mapsto AX-XB$ preserving $V$ reduces on $V$ to
$\operatorname{ad}_{A_0}+c\mathbf1$, and preservation of the complex volume
form forces $c=0$.  Hence the minimal bimodule supplies only the three
spin-1/quaternion-adjoint directions; the five quadrupole directions required
for full $su(3)$ cannot be supplied by a single left/right multiplication
pair.  We also record the independent Maurer--Cartan obstruction to obtaining
generic curvature from a pure local frame.  Finally, conditional on a genuine
$End_\C(V)$ colour connection entering a gauge-fixed Laplace-type
$\Theta$ Hessian, the standard heat-kernel $a_4$ coefficient necessarily
contains a Yang--Mills curvature-squared term.  Thus the remaining UBT problem
is sharply localized: derive the non-pure-gauge $End_\C(V)$ connection and its
normalisation from the finalized single action, rather than inserting it.
\end{abstract}

\section{Setup}

Let
\[
 \Aalg=\C\otimes_{\mathbb R}\Hq\simeq M_2(\C)
 =\C\mathbf1\oplus V,
 \qquad
 V=\C\operatorname{-span}\{I,J,K\}.
\]
Equivalently, in the matrix realisation $V$ is the traceless subspace
$sl_2(\C)$ as a complex vector space.  The canonical tensors on $V$ are
\[
 h(v,w)=\Sc(v^\dagger w),\qquad
 \Omega(v,w,u)=-\Sc(vwu),
\]
with joint stabiliser $SU(3)$.

The minimal two-sided operator used throughout the biquaternionic connection
ansatz is
\[
 T_{A,B}(X)=AX-XB,\qquad A,B\in M_2(\C).
\]

\section{Exact minimal-bimodule no-go}

\begin{theorem}[Restriction of a $V$-preserving bimodule operator]
\label{thm:restriction}
If $T_{A,B}(V)\subseteq V$, then there exists $c\in\C$ such that
\[
 A-B=c\mathbf1.
\]
Writing $A=a_0\mathbf1+A_0$ with $\operatorname{tr}A_0=0$, the restriction to
$V$ is therefore
\[
 \boxed{T_{A,B}|_V=\operatorname{ad}_{A_0}+c\mathbf1_V.}
\]
Consequently the complex vector space of such restrictions has dimension at
most four: three adjoint $sl_2(\C)$ directions plus one scalar direction.
\end{theorem}

\begin{proof}
For every traceless $X$,
\[
 0=\operatorname{tr}(T_{A,B}X)
 =\operatorname{tr}((A-B)X).
\]
The annihilator of the traceless matrices under the nondegenerate trace
pairing is exactly $\C\mathbf1$, hence $A-B=c\mathbf1$.  Substitution gives
\[
 AX-XB=AX-X(A-c\mathbf1)=[A,X]+cX=[A_0,X]+cX.
\]
The traceless $A_0$ space is three-dimensional over $\C$.
\end{proof}

\begin{corollary}[Volume preservation removes the scalar direction]
\label{cor:volume}
If in addition the infinitesimal action preserves the nonzero complex volume
form $\Omega\in\Lambda^3V^*$, then $c=0$.  Hence
\[
 \boxed{T|_V=\operatorname{ad}_{A_0}}
\]
and the available complex Lie algebra is only the three-dimensional
$\operatorname{ad}(sl_2(\C))$.
\end{corollary}

\begin{proof}
An infinitesimal endomorphism preserves a top-degree volume form iff its trace
on $V$ vanishes.  Since
$\operatorname{tr}_V(\operatorname{ad}_{A_0})=0$ and
$\operatorname{tr}_V(c\mathbf1_V)=3c$, volume preservation implies $c=0$.
\end{proof}

\begin{corollary}[No full $su(3)$ from one minimal left/right pair]
The minimal $A_\mu\Theta-\Theta B_\mu$ connection cannot realise the full
complexified colour algebra $sl_3(\C)$ on the canonical carrier while
preserving $V$ and $\Omega$.  It supplies precisely the quaternion-adjoint
spin triplet.  The five additional symmetric-traceless quadrupole generators
identified in the companion spin--quadrupole construction lie outside the
single-pair minimal bimodule image.
\end{corollary}

\begin{remark}
The statement is stronger than a raw parameter count.  Without imposing
$V$-preservation, the image of $(A,B)\mapsto L_A-R_B$ has complex dimension
seven because the common central pair $(c\mathbf1,c\mathbf1)$ is its kernel.
After imposing preservation of the canonical colour carrier the restriction
collapses to four complex directions, and $\Omega$-preservation leaves three.
\end{remark}

\section{Global stabiliser does not by itself imply a local gauge field}

\begin{proposition}[Fixed-carrier global-to-local obstruction]
Let $V$, $h$ and $\Omega$ be fixed over spacetime and let a kinetic term use
the ordinary derivative of a $V$-valued field.  The joint stabiliser
$SU(3)=\operatorname{Stab}(h,\Omega)$ is then a global symmetry.  Replacing a
constant $U$ by $U(x)$ creates the inhomogeneous term
$(\partial_\mu U)\phi$; local covariance requires an additional connection or
an equivalent nontrivial composite bundle construction.
\end{proposition}

\begin{proof}
For $\phi\mapsto U(x)\phi$,
\[
 \partial_\mu(U\phi)
 =U\partial_\mu\phi+(\partial_\mu U)\phi.
\]
The second term cannot be cancelled by the fixed tensors $h$ and $\Omega$.
\end{proof}

\section{Pure-frame Maurer--Cartan no-go}

\begin{theorem}[A pure local frame is flat]
If a candidate connection is constructed only as
\[
 \mathcal G=U^{-1}dU
\]
from a smooth local $SU(3)$ frame, then
\[
 \boxed{\mathcal F=d\mathcal G+\mathcal G\wedge\mathcal G=0}
\]
identically on that patch.
\end{theorem}

\begin{proof}
This is the Maurer--Cartan identity.
\end{proof}

Thus a changing basis of a fixed colour space cannot by itself provide a
generic gluon field.  A nonzero composite curvature would require additional
geometry, for example a genuinely varying subbundle/projector, or a
connection selected by the action rather than a pure frame.

\section{Minimal operator completion}

The already verified quaternion-adjoint generators
\[
 S_i=\frac{i}{2}\operatorname{ad}_{E_i}
\]
supply the three spin directions.  Define the symmetric traceless quadratic
operators
\[
 Q_{ij}=\frac12\{S_i,S_j\}-\frac23\delta_{ij}\mathbf1_V.
\]
Five independent $Q_{ij}$ together with the three $S_i$ span all eight
Gell--Mann directions.  Therefore the smallest presently known
biquaternionically generated \emph{operator} carrier for a colour connection
is
\[
 \mathfrak g_{\rm op}
 =\operatorname{span}_{\mathbb R}\{S_i,Q_{ij}^{\rm ST}\}\simeq su(3),
\]
and a candidate local connection has the form
\[
 \mathcal G_\mu
 =a_\mu^iS_i+b_\mu^{ij}Q_{ij}^{\rm ST}.
\]
This formula is kinematic.  The coefficient fields $a_\mu^i,b_\mu^{ij}$ are
not yet derived from $S_{\rm UBT}[\Theta]$.

\section{Conditional induced Yang--Mills term}

Suppose the finalized, gauge-fixed Euclidean $\Theta$ Hessian on a physical
fluctuation bundle is of Laplace type
\[
 P=-\left(g^{\mu\nu}\nabla_\mu\nabla_\nu+\mathcal E\right),
\]
and suppose $\nabla_\mu$ contains the colour connection $\mathcal G_\mu$.
Let
\[
 \mathcal F_{\mu\nu}=[\nabla_\mu,\nabla_\nu]_{\rm colour}.
\]
The local four-dimensional heat-kernel coefficient contains
\[
 \boxed{
 a_4(P)\supset
 \frac{1}{(4\pi)^2}\int d^4x\sqrt g\,
 \frac1{12}\operatorname{tr}
 \left(\mathcal F_{\mu\nu}\mathcal F^{\mu\nu}\right)
 }
\]
up to the stated Laplace/sign conventions and terms involving spacetime
curvature and $\mathcal E$.

For a real bosonic determinant
\[
 \Gamma^{(1)}
 =-\frac12\int_{\Lambda^{-2}}^\infty\frac{ds}{s}
 \operatorname{Tr}e^{-sP},
\]
this produces a logarithmic curvature-squared contribution.  If the physical
spectrum contains a compact-$\psi$ tower with masses
$m_n^2=m^2+n^2/R_\psi^2$, the common proper-time factor becomes
\[
 \mathcal I_0(m,R_\psi,\Lambda)
 =\int_{\Lambda^{-2}}^\infty\frac{ds}{s}\,
 e^{-m^2s}\,
 \vartheta_3(0,e^{-s/R_\psi^2}).
\]
Thus the same conditional spectral branch that yields an Einstein--Hilbert
term from the lower heat-kernel coefficient also has a standard route to an
induced Yang--Mills term at the next coefficient, provided the colour
connection is genuinely present in the Hessian.

\begin{tcolorbox}[colback=yellow!8!white,colframe=orange!70!black,title=Exact status]
\textbf{Closed [L1] / machine checked:}
the minimal-bimodule restriction theorem and its volume-preserving corollary.

\medskip
\textbf{Closed as a no-go:}
a single $V$-preserving, $\Omega$-preserving left/right pair does not furnish
full $su(3)$; a pure Maurer--Cartan frame has zero local curvature.

\medskip
\textbf{Closed conditionally:}
if a genuine colour connection occurs in a Laplace-type physical Hessian, the
standard $a_4$ coefficient contains $\operatorname{tr}\mathcal F^2$.

\medskip
\textbf{Open:}
derive $\mathcal G_\mu$, its eight coefficient fields, physical mode quotient
and normalisation from the finalized single UBT action.  The current action
registry remains \texttt{DEFINED\_FAMILY\_NOT\_FINALIZED}.
\end{tcolorbox}

\end{document}
